%% GENERATED by tools/from_docs.py from stepss-docs. Do not edit by hand:
%% edit the page named below in stepss-docs and regenerate.
%% source: stepss-docs/src/content/docs/developer/codegen-examples.mdx

\chapter{Worked CODEGEN examples}\label{doc:codegen-examples}

This page provides complete, annotated CODEGEN model files for each of the four model categories. These are real models from the RAMSES model library and can serve as templates for developing your own models.

For the block reference, see \hyperref[ch-libray_blocks]{CODEGEN Blocks}. For the model framework (states, equations, initialization), see \hyperref[ch-user_models]{User-Defined Models}.

\begin{center}\rule{0.5\linewidth}{0.5pt}\end{center}

\section{\texorpdfstring{Exciter Model: \texttt{exc\_ST1A}}{Exciter Model: exc\_ST1A}}\label{codegen-examples:exciter-model-exc_st1a}

The IEEE ST1A is a bus-fed static excitation system. This model demonstrates the typical structure of an \texttt{exc}-type CODEGEN file: voltage compensation, AVR with lead-lag stages, field current limiting, UEL/OEL gating, and variable-limit output clamp.

\textbf{Source:} IEEE\_models/exc\_ST1A.txt (from the RAMSES model library)

\textbf{Full Model File}

\begin{Verbatim}
exc
ST1A

%data

Kv
Rc
Xc
TR
UEL
VIMIN
VIMAX
VUEL    ERROR - Vuel is the UEL signal not a parameter!
TC
TB
TC1
TB1
KA
TA
VAMIN
VAMAX
VRMIN
VRMAX
KC
KF
TF
KLR
ILR

%parameters

VREF     = vcomp([v],[p],[q],{Kv},{Rc},{Xc})+([vf]/{KA})

%states

Vc1      = vcomp([v],[p],[q],{Kv},{Rc},{Xc})
Vc       = vcomp([v],[p],[q],{Kv},{Rc},{Xc})
deltaV   = ([vf]/{KA})
V1       = [vf]/{KA}
V2       = [vf]/{KA}
V3       = [vf]/{KA}
V4       = [vf]/{KA}
VA       = [vf]
VLR      = {KLR}*([if]-{ILR})
VLRlim   = 0.
V5       = [vf]
V6       = [vf]
V7       = [vf]
uplim    = [v]*{VRMAX}-{KC}*[if]
lolim    = [v]*{VRMIN}
VF       = 0.

VOEL     = 99999.

%observables

VA
Vc1
Vc
V1
V2
V3
V4
vf
p
q
v
VLRlim
deltaV

VREF
VOEL

%models
& algeq                    fictitious over excitation limit
[VOEL]-99999.
& algeq                    line drop compensation
[Vc1]-vcomp([v],[p],[q],{Kv},{Rc},{Xc})
& tf1p                     voltage measurement time constant
Vc1
Vc
1.d0
{TR}
& algeq                    summation point of AVR
[deltaV]-{VREF}+[Vc]-equal({UEL},1.d0)*{VUEL}+[VF]
& lim                      limiter on deltaV
deltaV
V1
{VIMIN}
{VIMAX}
& max1v1c                  HV gate for VUEL if UEL=2
V1
V2
{VUEL}*equal({UEL},2.d0) - 9999.*(1-equal({UEL},2.d0))
& tf1p1z                   first lead-lag
V2
V3
1.
{TC}
{TB}
& tf1p1z                   second lead-lag
V3
V4
1.
{TC1}
{TB1}
& tf1plim                  amplifier
V4
VA
{KA}
{TA}
{VAMIN}
{VAMAX}
& algeq                    immediate field current limiting signal
[VLR]-{KLR}*([if]-{ILR})
& lim                      ... lower limited to zero
VLR
VLRlim
0.
99999.
& algeq                    ... and subtracted from main signal
[V5]-[VA]+[VLRlim]
& max1v1c                  HV gate for VUEL if UEL=3
V5
V6
{VUEL}*equal({UEL},3.d0) - 9999.*(1-equal({UEL},3.d0))
& min2v                    LV gate for VOEL
V6
VOEL
V7
& algeq                    variable lower limit of final limiter
[v]*{VRMIN}-[lolim]
& algeq                    variable upper limit of final limiter
[v]*{VRMAX}-{KC}*[if]-[uplim]
& limvb                    final limiter
V7
lolim
uplim
vf
& tfder1p                  derivative feedback
V7
VF
{KF}/{TF}
{TF}
\end{Verbatim}

\textbf{Structure Guide}

\subsection{File Header}\label{codegen-examples:file-header}

\begin{Verbatim}
exc          <- model type (exciter)
ST1A         <- model name (used in the dynamic data file as EXC ST1A)
\end{Verbatim}

\subsection{\texorpdfstring{\texttt{\%data} Section}{\%data Section}}\label{codegen-examples:data-section}

Lists the parameters that the user provides in the dynamic data file. Each name becomes a \texttt{\{name\}} reference in expressions. The order here defines the order of values in the data file.

\subsection{\texorpdfstring{\texttt{\%parameters} Section}{\%parameters Section}}\label{codegen-examples:parameters-section}

Operating-point parameters computed at initialization. The number of parameters must equal the number of output states (here 1: \texttt{vf}). \texttt{VREF} is the voltage setpoint, computed from the initial operating point so that the steady-state field voltage is matched.

\subsection{\texorpdfstring{\texttt{\%states} Section}{\%states Section}}\label{codegen-examples:states-section}

All internal and output states with their initial values. The initialization expressions use system variables (\texttt{{[}v{]}}, \texttt{{[}p{]}}, \texttt{{[}q{]}}, \texttt{{[}if{]}}, \texttt{{[}vf{]}}) and data values (\texttt{\{KA\}}, etc.). The last state \texttt{vf} is the output state (field voltage).

\subsection{\texorpdfstring{\texttt{\%observables} Section}{\%observables Section}}\label{codegen-examples:observables-section}

States that can be recorded during simulation.

\subsection{\texorpdfstring{\texttt{\%models} Section: Signal Flow}{\%models Section: Signal Flow}}\label{codegen-examples:models-section-signal-flow}

\begin{enumerate}
\def\labelenumi{\arabic{enumi}.}
\tightlist
\item
  \texttt{algeq}, OEL placeholder (set to large value, overridden by OEL model if present)
\item
  \texttt{algeq}, line drop compensation: \(V_{c1} = \text{vcomp}(V, P, Q, K_v, R_c, X_c)\)
\item
  \texttt{tf1p}, voltage measurement filter with time constant \(T_R\)
\item
  \texttt{algeq}, AVR summation: \(\Delta V = V_{REF} - V_c - V_{UEL} + V_F\)
\item
  \texttt{lim}, error limiter \([V_{IMIN}, V_{IMAX}]\)
\item
  \texttt{max1v1c}, HV gate for UEL (active when UEL=2)
\item
  \texttt{tf1p1z}, first lead-lag: \((1 + sT_C)/(1 + sT_B)\)
\item
  \texttt{tf1p1z}, second lead-lag: \((1 + sT_{C1})/(1 + sT_{B1})\)
\item
  \texttt{tf1plim}, amplifier with gain \(K_A\), time constant \(T_A\), limits \([V_{AMIN}, V_{AMAX}]\)
\item
  \texttt{algeq} + \texttt{lim}, instantaneous field current limiter \(V_{LR} = K_{LR}(I_f - I_{LR})\), lower-limited to 0
\item
  \texttt{algeq}, subtract limiter signal from main path
\item
  \texttt{max1v1c}, HV gate for UEL (active when UEL=3)
\item
  \texttt{min2v}, LV gate for OEL
\item
  \texttt{algeq} \ensuremath{\times} 2, compute variable limits: \(V_{RMIN} \cdot V_t\) and \(V_{RMAX} \cdot V_t - K_C I_f\)
\item
  \texttt{limvb}, bus-fed variable-limit output clamp \ensuremath{\rightarrow} \texttt{vf}
\item
  \texttt{tfder1p}, derivative feedback: \(sK_F/T_F / (1 + sT_F)\)
\end{enumerate}

\begin{center}\rule{0.5\linewidth}{0.5pt}\end{center}

\section{\texorpdfstring{Governor Model: \texttt{tor\_hygov}}{Governor Model: tor\_hygov}}\label{codegen-examples:governor-model-tor_hygov}

A hydraulic turbine-governor model with speed droop, PI controller, rate limiter, gate servo, and non-linear turbine (head-flow-power). This demonstrates a \texttt{tor}-type model with both differential and algebraic equations.

\textbf{Source:} IEEE\_models/tor\_hygov.txt (from the RAMSES model library)

\textbf{Full Model File}

\begin{Verbatim}
tor
hygov

%data

R
r
Tr
Tf
Tg
VELM
Gmax
Gmin
Tw
At
Dturb
Qnl

%parameters

Po = ([p]/{At})+{Qnl}    ! power setpoint

%states

Q = {Po}    ! water flow
dH = 0.     ! deviation of head from 1.
g = {Po}    ! gate opening
c = {Po}    ! desired gate opening
x5 = {Po}   ! output of PI controller
x4 = {Po}   ! output of integrator of PI controller
x3 = 0.     ! output of rate limiter = input of integrator of PI controller
x2 = 0.     ! input of rate limiter of PI controller
e = 0.      ! output of Tf time constant block
x1 = 0.     ! input of Tf time constant block
Pm = [p]    ! mechanical power

%observables

Pm
Q
dH
g

%models

& algeq    ! input of Tf time constant block
1.-[omega]-{R}*([c]-{Po}) - [x1]
& tf1p     ! Tf time constant
x1
e
1.
{Tf}
& algeq    ! input of rate limiter
[e]/({r}*{Tr}) - [x2]
& lim      ! rate limiter
x2
x3
-{VELM}
 {VELM}
& inlim    ! integrator of PI controller
x3
x4
1.
{Gmin}
{Gmax}
& algeq    ! output of PI controller
[x4]+([e]/{r}) - [x5]
& lim      ! gate limiter
x5
c
{Gmin}
{Gmax}
& tf1p     ! speed governor : time constant
c
g
1.
{Tg}

& algeq                 ! hydro turbine: head as function of water flow and gate opening
1. -([Q]/[g])**2 -[dH]
& int                   ! hydro turbine: water column inertia
dH
Q
{Tw}
& algeq                 ! hydro turbine: mechanical torque
{At}*(1.-[dH])*([Q]-{Qnl})-{Dturb}*(1.-[omega])*[g]  - [tm]*[omega]
& algeq
[tm]*[omega] - [Pm]
\end{Verbatim}

\textbf{Structure Guide}

\subsection{File Header}\label{codegen-examples:file-header-1}

\begin{Verbatim}
tor          <- model type (torque/governor)
hygov        <- model name (used as `TOR HYGOV` in the data file; dispatcher case is `tor_HYGOV`)
\end{Verbatim}

\subsection{Input/Output States}\label{codegen-examples:inputoutput-states}

\begin{itemize}
\tightlist
\item
  \textbf{Input:} \texttt{{[}omega{]}} (rotor speed), \texttt{{[}p{]}} (electrical power)
\item
  \textbf{Output:} \texttt{{[}tm{]}} (mechanical torque)
\end{itemize}

\subsection{Signal Flow}\label{codegen-examples:signal-flow}

\textbf{Governor section:}
1. Speed error with droop: \(x_1 = 1 - \omega - R \cdot (g - P_0)\)
2. Filter: \texttt{tf1p} with time constant \(T_f\)
3. PI controller path: rate limiter (\texttt{lim}) \ensuremath{\rightarrow} integrator (\texttt{inlim}) with gate limits \ensuremath{\rightarrow} proportional bypass
4. Gate servo: \texttt{tf1p} with time constant \(T_g\)

\textbf{Turbine section:}
5. Head-flow relation: \(\Delta H = 1 - (Q/g)^2\) (non-linear algebraic)
6. Water column inertia: \(T_w \dot{Q} = \Delta H\) (\texttt{int} block)
7. Mechanical torque: \(T_m \omega = A_t (1 - \Delta H)(Q - Q_{nl}) - D_{turb}(1 - \omega) g\)
8. Mechanical power: \(P_m = T_m \omega\)

\subsection{Key Parameters}\label{codegen-examples:key-parameters}

{\def\LTcaptype{none} % do not increment counter
\begin{small}\begin{longtable}[]{@{}ll@{}}
\toprule\noalign{}
Parameter & Description \\
\midrule\noalign{}
\endhead
\bottomrule\noalign{}
\endlastfoot
\texttt{R} & Permanent droop \\
\texttt{r} & Temporary droop \\
\texttt{Tr} & Reset time \\
\texttt{Tf} & Filter time constant \\
\texttt{Tg} & Gate servo time constant \\
\texttt{VELM} & Gate velocity limit \\
\texttt{Gmax/Gmin} & Gate position limits \\
\texttt{Tw} & Water starting time \\
\texttt{At} & Turbine gain \\
\texttt{Dturb} & Turbine damping \\
\texttt{Qnl} & No-load flow \\
\end{longtable}\end{small}
}

\begin{center}\rule{0.5\linewidth}{0.5pt}\end{center}

\section{\texorpdfstring{Injector Model: \texttt{inj\_vfd\_load}}{Injector Model: inj\_vfd\_load}}\label{codegen-examples:injector-model-inj_vfd_load}

A voltage- and frequency-dependent load model with exponential characteristics and low-voltage protection (constant admittance below a threshold). This is the most widely used load model in RAMSES. It demonstrates an \texttt{inj}-type model with current injection equations.

\textbf{Source:} custom\_models/inj\_vfd\_load.txt (from the RAMSES model library)

\textbf{Full Model File}

\begin{Verbatim}
inj
vfd_load

%data

Dp         ! sensitivity of active power to frequency
a1         ! proportion of type 1 in active power
alpha1     ! exponent of type 1 in active power
a2         ! proportion of type 2 in active power
alpha2     ! exponent of type 2 in active power
alpha3     ! exponent of type 3 in active power  (proportion = 1-a1-a2)
Dq         ! sensitivity of reactive power to frequency
b1         ! proportion of type 1 in reactive power
beta1      ! exponent of type 1 in reactive power
b2         ! proportion of type 2 in reactive power
beta2      ! exponent of type 2 in reactive power
beta3      ! exponent of type 3 in reactive power  (proportion = 1-b1-b2)
Vinit      ! initial voltage at equivalent distribution bus
Vlow       ! voltage below which the load behaves as constant admittance
foption    ! set to 1 to use bus frequency; otherwise COI speed
Tmes       ! measurement time constant for bus frequency estimation

%parameters

Vo = equal({Vinit},0.d0)*dsqrt([vx]**2+[vy]**2) + {Vinit}
r =  dsqrt([vx]**2+[vy]**2)/{Vo}
denomP = {a1}*{Vo}**{alpha1}+{a2}*{Vo}**{alpha2}+(1.d0-{a1}-{a2})*{Vo}**{alpha3}
denomQ = {b1}*{Vo}**{beta1}+{b2}*{Vo}**{beta2}+(1.d0-{b1}-{b2})*{Vo}**{beta3}
Po = -[vx]*[ix]-[vy]*[iy]
Qo =  [vx]*[iy]-[vy]*[ix]
Vmin = max(1.d-01,{Vlow})
Geq = {Po}*({a1}*{Vmin}**{alpha1}+{a2}*{Vmin}**{alpha2}
      +(1.d0-{a1}-{a2})*{Vmin}**{alpha3})/({denomP}*{Vmin}**2)
Beq =-{Qo}*({b1}*{Vmin}**{beta1}+{b2}*{Vmin}**{beta2}
      +(1.d0-{b1}-{b2})*{Vmin}**{beta3})/({denomQ}*{Vmin}**2)

%states

Vreal = {Vo}      ! voltage at equivalent distribution bus
V = {Vo}          ! Vreal slightly delayed
P = {Po}          ! active power
Q = {Qo}          ! reactive power
fbus = 1.d0       ! bus frequency
df = 0.d0         ! frequency deviation
u = 1             ! model switch: =1 if V > Vmin, =0 otherwise

%observables

P
Q
df
u

%models

& f_inj
fbus
{Tmes}
& algeq
equal(1.d0,{foption})*([fbus]-1.d0)
  + (1.d0-equal(1.d0,{foption}))*([omega]-1.d0) - [df]
& algeq         ! voltage at distribution bus = network voltage / trfo ratio
dsqrt([vx]**2+[vy]**2)/{r}-[Vreal]
& tf1p          ! small measurement time constant
Vreal
V
1.
0.003
& pwlin4        ! determine if u=1 or u=0
V
u
-1.d0
0.d0
{Vmin}
0.d0
{Vmin}
1.d0
2.d0
1.d0
& algeq         ! active power function of V, u and f
[u]*({Po}*(1.d0+{Dp}*[df])
  *({a1}*[V]**{alpha1}+{a2}*[V]**{alpha2}
    +(1.d0-{a1}-{a2})*[V]**{alpha3})/{denomP})
  + (1.d0-[u])*{Geq}*[V]**2 -[P]
& algeq         ! reactive power function of V, u and f
[u]*({Qo}*(1.d0+{Dq}*[df])
  *({b1}*[V]**{beta1}+{b2}*[V]**{beta2}
    +(1.d0-{b1}-{b2})*[V]**{beta3})/{denomQ})
  - (1.d0-[u])*{Beq}*[V]**2 -[Q]
& algeq         ! ix component of injected current
(-[P]*[vx] - [Q]*[vy])/max([V]**2,1.d-02) -[ix]
& algeq         ! iy component of injected current
( [Q]*[vx] - [P]*[vy])/max([V]**2,1.d-02) -[iy]
\end{Verbatim}

\textbf{Structure Guide}

\subsection{File Header}\label{codegen-examples:file-header-2}

\begin{Verbatim}
inj          <- model type (injector)
vfd_load     <- model name (used as INJEC vfd_load in the data file)
\end{Verbatim}

\subsection{Input/Output States}\label{codegen-examples:inputoutput-states-1}

\begin{itemize}
\tightlist
\item
  \textbf{Input:} \texttt{{[}vx{]}}, \texttt{{[}vy{]}} (bus voltage components), \texttt{{[}omega{]}} (COI frequency)
\item
  \textbf{Output:} \texttt{{[}ix{]}}, \texttt{{[}iy{]}} (injected current components)
\end{itemize}

\subsection{Model Logic}\label{codegen-examples:model-logic}

\textbf{Frequency estimation:}
1. \texttt{f\_inj}, estimates bus frequency from voltage phasor derivatives
2. \texttt{algeq}, selects bus frequency or COI frequency based on \texttt{foption}

\textbf{Voltage processing:}
3. \texttt{algeq}, compute distribution bus voltage: \(V_{real} = \sqrt{v_x^2 + v_y^2}/r\)
4. \texttt{tf1p}, small delay (3 ms) for numerical stability
5. \texttt{pwlin4}, switch: \(u = 1\) if \(V > V_{min}\), \(u = 0\) otherwise

\textbf{Power computation:}
6. Active power: \(P = u \cdot P_0 (1 + D_p \Delta f) \frac{a_1 V^{\alpha_1} + a_2 V^{\alpha_2} + (1-a_1-a_2) V^{\alpha_3}}{\text{denom}_P} + (1-u) G_{eq} V^2\)
7. Reactive power: analogous with \(b\)/\(\beta\) coefficients

\textbf{Current injection:}
8--9. Convert P, Q to rectangular currents: \(i_x = (-PV_x - QV_y)/V^2\), \(i_y = (QV_x - PV_y)/V^2\)

\subsection{Key Design Patterns}\label{codegen-examples:key-design-patterns}

\begin{itemize}
\tightlist
\item
  \textbf{Low-voltage protection}: below \(V_{min}\), the load switches to constant admittance (\(G_{eq}\), \(B_{eq}\)) to prevent numerical issues
\item
  \textbf{Distribution transformer}: the ratio \(r\) models the voltage step-down to the load bus
\item
  \textbf{Frequency dependence}: \(D_p\), \(D_q\) scale power with frequency deviation
\end{itemize}

\begin{center}\rule{0.5\linewidth}{0.5pt}\end{center}

\section{\texorpdfstring{Two-Port Model: \texttt{twop\_HVDC\_LCC}}{Two-Port Model: twop\_HVDC\_LCC}}\label{codegen-examples:two-port-model-twop_hvdc_lcc}

A Line-Commutated Converter (LCC) HVDC link model connecting two AC buses through a DC line. This demonstrates a \texttt{twop}-type model with complex AC/DC interaction, multiple control modes (current control, voltage control, gamma control), and non-linear commutation equations.

\textbf{Source:} custom\_models/twop\_HVDC\_LCC.txt (from the RAMSES model library)

\textbf{Full Model File}

\begin{Verbatim}
twop
HVDC_LCC

%data

p          ! number of poles (1 or 2)
rec_ctl    ! rectifier control: 1=current, 2=power, -1=alpha min, 3=blocked
inv_ctl    ! inverter control: 1=voltage, 2=gamma, -2=gamma min, -1=reduced current, 3=blocked
Br         ! number of bridges in series in rectifier
Bi         ! number of bridges in series in inverter
Er         ! phase-to-phase open-circuit voltage (kV) on DC side of rectifier transformer
Ei         ! phase-to-phase open-circuit voltage (kV) on DC side of inverter transformer
Rcr        ! rectifier commutating resistance per bridge (ohm)
Xcr        ! rectifier commutating reactance per bridge (ohm)
Rci        ! inverter commutating resistance per bridge (ohm)
Xci        ! inverter commutating reactance per bridge (ohm)
RL         ! DC line series resistance (ohm)
Rcomp      ! compounding resistance for inverter voltage control (ohm)
Tp         ! time constant of power control (s)
Idset      ! initial DC current setpoint (kA)
nr         ! initial rectifier transformer ratio (pu/pu)
ni         ! initial inverter transformer ratio (pu/pu)

%parameters

Pset = -([vx1]*[ix1]+[vy1]*[iy1])*sbase1
Vdr0 = {Pset}/({p}*{Idset})
alpha0 = dacos( ({nr}/(1.3504745*{Er}*dsqrt([vx1]**2+[vy1]**2)))
         * (({Vdr0}/{Br})+(0.9549297*{Xcr}+2*{Rcr})*{Idset}) )
mur0 = dacos( dcos({alpha0})
       - 1.4142136*{Xcr}*{nr}*{Idset}
         /(dsqrt([vx1]**2+[vy1]**2)*{Er})) - {alpha0}
Qhvdc1 = {Pset}*(2.d0*{mur0}+dsin(2*{alpha0})-dsin(2d0*({mur0}+{alpha0})))
         /(dcos(2.d0*{alpha0})-dcos(2d0*({mur0}+{alpha0})))
B1 = ({Qhvdc1}/sbase1-[vx1]*[iy1]+[vy1]*[ix1])/([vx1]**2+[vy1]**2)
Vdi0 = {Vdr0}-{RL}*{Idset}
gamma0 = dacos( ({ni}/(1.3504745*{Ei}*dsqrt([vx2]**2+[vy2]**2)))
          * (({Vdi0}/{Bi})+(0.9549297*{Xci}+2*{Rci})*{Idset}) )
mui0 = dacos( dcos({gamma0})
       - 1.4142136*{Xci}*{ni}*{Idset}
         /(dsqrt([vx2]**2+[vy2]**2)*{Ei})) - {gamma0}
P2in = {p}*{Vdi0}*{Idset}
G2 = ({P2in}/sbase2-[vx2]*[ix2]-[vy2]*[iy2])
     /dsqrt([vx2]**2+[vy2]**2)
Qhvdc2 = {P2in}*(2.d0*{mui0}+dsin(2*{gamma0})-dsin(2d0*({mui0}+{gamma0})))
         /(dcos(2.d0*{gamma0})-dcos(2d0*({mui0}+{gamma0})))
B2 = ({Qhvdc2}/sbase2-[vx2]*[iy2]+[vy2]*[ix2])/([vx2]**2+[vy2]**2)
Vset = {Vdi0}+{Rcomp}*{Idset}
alphamin = 0.d0
Idred = 0.d0
gammamin = 0.d0
Td = 0.02

%states

Vdr = {Vdr0}         ! rectifier DC voltage (kV)
alpha = {alpha0}      ! firing angle (rad)
mur = {mur0}          ! overlap angle (rad)
Id = {Idset}          ! DC current (kA)
Vdi = {Vdi0}          ! inverter DC voltage (kV)
gamma = {gamma0}      ! extinction angle (rad)
mui = {mui0}          ! inverter overlap angle (rad)
Ides = {Idset}        ! desired DC current (kA)
Iord = {Idset}        ! DC current order (kA)
P1 = (-[vx1]*[ix1]-[vy1]*[iy1])*sbase1
Q1 = ([vx1]*[iy1]-[vy1]*[ix1])*sbase1
P2 = ([vx2]*[ix2]+[vy2]*[iy2])*sbase2
Q2 = ([vx2]*[iy2]-[vy2]*[ix2])*sbase2
V1 = dsqrt([vx1]**2+[vy1]**2)
V1d = dsqrt([vx1]**2+[vy1]**2)
V2 = dsqrt([vx2]**2+[vy2]**2)
V2d = dsqrt([vx2]**2+[vy2]**2)
alphad = {alpha0}*57.2957795
gammad = {gamma0}*57.2957795

%observables

Vdr
alphad
mur
Vdi
gammad
mui
Ides
Iord
Id
P1
Q1
P2
Q2
Pset

%models

& algeq          ! Vdr: rectifier DC voltage equation
[Vdr] - {Br}*( 1.3504745*[V1d]*{Er}*dcos([alpha])/{nr}
  - (0.9549297*{Xcr}+2*{Rcr})*[Id] )
& algeq          ! overlap angle at rectifier
dcos([alpha]+[mur]) - dcos([alpha])
  + 1.4142136*{Xcr}*{nr}*[Id]/(max([V1d],1.d-03)*{Er})
& algeq          ! ix1: rectifier AC current (x-component)
{B1}*[vy1]-([vx1]+[vy1]*(2*[mur]+dsin(2*[alpha])
  -dsin(2*([mur]+[alpha])))/max((dcos(2*[alpha])
  -dcos(2*([mur]+[alpha]))),1d-03))
  *{p}*[Vdr]*[Id]/(sbase1*max([V1d]**2,1d-06)) - [ix1]
& algeq          ! iy1: rectifier AC current (y-component)
-{B1}*[vx1]-([vy1]-[vx1]*(2*[mur]+dsin(2*[alpha])
  -dsin(2*([mur]+[alpha])))/max((dcos(2*[alpha])
  -dcos(2*([mur]+[alpha]))),1d-03))
  *{p}*[Vdr]*[Id]/(sbase1*max([V1d]**2,1d-06)) - [iy1]
& algeq          ! Vdi: inverter DC voltage equation
[Vdi] - {Bi}*( 1.3504745*[V2d]*{Ei}*dcos([gamma])/{ni}
  - (0.9549297*{Xci}+2*{Rci})*[Id] )
& algeq          ! overlap angle at inverter
dcos([gamma]+[mui]) - dcos([gamma])
  + 1.4142136*{Xci}*{ni}*[Id]/(max([V2d],1.d-03)*{Er})
& algeq          ! ix2: inverter AC current (x-component)
-{G2}*[vx2]+{B2}*[vy2]+([vx2]-[vy2]*(2*[mui]+dsin(2*[gamma])
  -dsin(2*([mui]+[gamma])))/max((dcos(2*[gamma])
  -dcos(2*([mui]+[gamma]))),1d-03))
  *{p}*[Vdi]*[Id]/(sbase2*max([V2d]**2,1d-06)) - [ix2]
& algeq          ! iy2: inverter AC current (y-component)
-{B2}*[vx2]-{G2}*[vy2]+([vy2]+[vx2]*(2*[mui]+dsin(2*[gamma])
  -dsin(2*([mui]+[gamma])))/max((dcos(2*[gamma])
  -dcos(2*([mui]+[gamma]))),1d-03))
  *{p}*[Vdi]*[Id]/(sbase2*max([V2d]**2,1d-06)) - [iy2]
& algeq          ! DC link: Kirchhoff voltage law
[Vdr]-[Vdi]-{RL}*[Id]

& algeq          ! desired DC current (depends on control mode)
(equal({rec_ctl},1.d0)*{Idset})
  +(equal({rec_ctl},2.d0)*({Pset}/max([Vdr],1d-02)))
  -(equal({rec_ctl},3.d0)*50.d0)
  -(equal({rec_ctl},-1.d0)*50.d0) -[Ides]
& tf1plim        ! current order (filtered and limited)
Ides
Iord
1.d0
Tp
0.d0
9999.d0

& algeq          ! rectifier control equation
equal({rec_ctl},1.d0)*([Iord]-[Id])
  +equal({rec_ctl},2.d0)*([Iord]-[Id])
  +equal({rec_ctl},-1d0)*({alphamin}-[alpha])
  +equal({rec_ctl},3.d0)*[Id]
& algeq          ! inverter control equation
equal({inv_ctl},1.d0)*([Vdi]+{Rcomp}*[Id]-{Vset})
  +equal({inv_ctl},2.d0)*({gamma0}-[gamma])
  +equal({inv_ctl},-2d0)*({gammamin}-[gamma])
  +equal({inv_ctl},-1d0)*({Idred}-[Id])
  +equal({inv_ctl},3.d0)*(1.570796327-[gamma])

& algeq          ! AC voltage at bus 1
dsqrt([vx1]**2+[vy1]**2) - [V1]
& tf1p           ! delayed AC voltage at bus 1
V1
V1d
1.
Td
& algeq          ! AC voltage at bus 2
dsqrt([vx2]**2+[vy2]**2) - [V2]
& tf1p           ! delayed AC voltage at bus 2
V2
V2d
1.
Td
& algeq          ! observable: P1
(-[vx1]*[ix1]-[vy1]*[iy1])*sbase1-[P1]
& algeq          ! observable: Q1
( [vx1]*[iy1]-[vy1]*[ix1])*sbase1-[Q1]
& algeq          ! observable: P2
([vx2]*[ix2]+[vy2]*[iy2])*sbase2-[P2]
& algeq          ! observable: Q2
([vx2]*[iy2]-[vy2]*[ix2])*sbase2-[Q2]
& algeq          ! alpha in degrees
[alphad]-[alpha]*57.2957795
& algeq          ! gamma in degrees
[gammad]-[gamma]*57.2957795
\end{Verbatim}

\textbf{Structure Guide}

\subsection{File Header}\label{codegen-examples:file-header-3}

\begin{Verbatim}
twop         <- model type (two-port)
HVDC_LCC     <- model name (used as `TWOP HVDC_LCC ...` in data file)
\end{Verbatim}

\subsection{Input/Output States}\label{codegen-examples:inputoutput-states-2}

\begin{itemize}
\tightlist
\item
  \textbf{Input:} \texttt{{[}vx1{]}}, \texttt{{[}vy1{]}}, \texttt{{[}vx2{]}}, \texttt{{[}vy2{]}} (voltage at both buses)
\item
  \textbf{Output:} \texttt{{[}ix1{]}}, \texttt{{[}iy1{]}}, \texttt{{[}ix2{]}}, \texttt{{[}iy2{]}} (current injected at both buses)
\item
  System variables: \texttt{sbase1}, \texttt{sbase2} (base power at each bus)
\end{itemize}

\subsection{Model Structure}\label{codegen-examples:model-structure}

\textbf{DC equations (algebraic):}
- Rectifier DC voltage: \(V_{dr} = B_r \left(\frac{3\sqrt{2}}{\pi} V_1 E_r \cos\alpha / n_r - (R_{cr} + \frac{3}{\pi}X_{cr}) I_d\right)\)
- Inverter DC voltage: analogous with \(\gamma\), \(E_i\), \(n_i\)
- Overlap angles: from commutation equations
- DC link: \(V_{dr} - V_{di} = R_L I_d\)

\textbf{AC current injection:}
- Includes fundamental-frequency reactive power from commutation overlap
- Shunt compensation (\(B_1\), \(B_2\), \(G_2\)) to match initial operating point

\textbf{Control modes} (selected by \texttt{rec\_ctl} and \texttt{inv\_ctl} using \texttt{equal()} function):
- Rectifier: constant current, constant power, minimum alpha, blocked
- Inverter: constant voltage, constant gamma, minimum gamma, reduced current, blocked

\subsection{Key Design Patterns}\label{codegen-examples:key-design-patterns-1}

\begin{itemize}
\tightlist
\item
  \textbf{\texttt{equal()} function}: used to switch between control modes without discrete variables, acts as a continuous selector
\item
  \textbf{\texttt{sbase1}/\texttt{sbase2}}: system variables for per-unit conversion at each bus
\item
  \textbf{\texttt{Td} delay}: small artificial time constant on terminal voltages for numerical robustness
\item
  \textbf{Initialization}: complex parameter expressions compute initial angles, overlap, and compensation from the power flow solution
\end{itemize}
